\documentclass[12pt,a4paper]{article}

\usepackage[utf8]{inputenc}
\usepackage[T1]{fontenc}
\usepackage{times}
\usepackage[margin=2.5cm]{geometry}
\usepackage{amsmath,amssymb}
\usepackage{graphicx}
\usepackage{booktabs}
\usepackage{enumitem}
\usepackage{hyperref}
\usepackage{xcolor}
\usepackage{caption}
\usepackage{float}
\usepackage{parskip}

\hypersetup{colorlinks=true,linkcolor=blue,citecolor=blue,urlcolor=blue}

\begin{document}

\begin{center}
{\Large\textbf{ASSIGNMENT 3}}
\end{center}

\noindent\textbf{Student:} Truong Tri Dung -- MITIU25208

% ============================================================
\section{Review 1: A Research for Segmentation of Brain Tumors Based on GAN Model}
% ============================================================

The paper proposes a Multi-scale 3D-GAN model for brain tumor segmentation on the BraTS 2021 dataset, with three main contributions: a Generator architecture based on a 3D Residual U-Net combined with an intermediate ``Bridge'' module that aggregates multi-scale features before feeding them to the Discriminator, an algorithm for dynamically adjusting the number of Discriminator training epochs based on the linear regression slope of the adversarial loss in order to balance the G--D training process, and a loss function combining Dice loss with MSE adversarial loss; experimental results show that the model achieves Dice scores of 82.3\% (ET), 88.6\% (WT), and 83.1\% (TC), outperforming the U-Net, V-Net, VoxelGAN, and Vox2Vox baselines. In terms of strengths, the Bridge module is a well-motivated idea as it enables the Discriminator to exploit both high-level contextual information and low-level details simultaneously, the dynamic adjustment algorithm for Discriminator epochs is a novel and practical contribution to a classical problem in GAN training, and the paper provides a fairly comprehensive evaluation on four metrics (Dice, Jaccard, Hausdorff95, ASSD) and compares against four strong baselines covering both CNN-based and GAN-based approaches. However, the paper still has several notable limitations: the experimental reliability is not fully convincing since the entire evaluation is conducted only on a validation set split from the training data, without cross-validation and without testing on the official BraTS 2021 server, which makes the results difficult to compare directly with other works on the same dataset; the paper lacks an ablation study for the two core contributions --- the Bridge module and Algorithm 1 --- so it is impossible to determine the individual contribution of each component; there exists an unexplained paradox in which the Multi-scale GAN achieves a higher Dice score than the ResUnet variant without GAN, yet its Hausdorff and ASSD values on the Enhancing Tumor class are actually worse (9.6mm vs.\ 6.4mm; 3.37 vs.\ 2.35), suggesting that the GAN component may improve volume overlap but degrade boundary quality --- a phenomenon that should have been analyzed carefully; the important hyperparameters in Algorithm 1 ($k$, $l$, threshold $\beta$) are not investigated, and there is no theoretical justification for applying linear regression to the adversarial loss; furthermore, the paper does not report the model's parameter count, training and inference time, nor any statistical significance tests for the comparisons. In terms of presentation, the main contribution of the Bridge module deserves its own dedicated illustration rather than being merged into Fig.~1. To improve the quality of the paper, the authors should add evaluation on the BraTS online server or at least conduct 5-fold cross-validation with mean $\pm$ std, perform a thorough ablation study for both the Bridge module and Algorithm 1, explain the Hausdorff/ASSD paradox on the ET class, investigate the sensitivity to the key hyperparameters, and release the source code to ensure reproducibility. Overall, the paper presents a reasonable idea and competitive Dice results, but requires substantial improvements in experimental rigor.

% ============================================================
\section{Review 2: Paradigm shift from ANNs to DCNNs in medical image processing}
% ============================================================

This paper is a review that discusses the shift from traditional Artificial Neural Networks (ANNs) to deep Convolutional Neural Networks (DCNNs) in the field of medical image processing, with the main argument being that in the past people had to manually extract features before feeding them into an ANN, whereas now DCNNs can learn directly from images without that step. I think the paper has several nice points: the explanation of the convolution operation using a concrete numerical matrix example and the application of a Sobel kernel to a hand X-ray image is very easy to understand and suitable for beginners; the two summary tables comparing ANN-based and DCNN-based studies are also quite convincing as they clearly show that the ``feature extraction'' column disappears entirely in DCNN-based studies; furthermore, the authors are upfront about the limitations of DCNNs such as the need for large datasets, the tendency to overfit, and the black-box nature of the models. However, I think the paper also has a number of shortcomings. First, although the paper was published in 2024, it makes no mention at all of Vision Transformers, SAM, or medical foundation models such as Med-PaLM or LLaVA-Med --- yet these are the directions that have been reshaping the field recently, and if we are going to talk about a ``paradigm shift'', then arguably the next shift is already happening, rather than the ANN$\rightarrow$DCNN transition which has already been widely acknowledged since around 2017. Second, the paper mainly revolves around the classification task and largely overlooks segmentation --- an important branch with models such as U-Net and nnU-Net --- and it also does not mention standard datasets such as BraTS, LIDC, or CheXpert. Third, the entire paper focuses on the English-language context and does not devote a single line to low-resource languages; I find this particularly unfortunate because I am currently working on a research direction related to clinical decision support systems for Vietnamese, so I understand that when applying DCNNs to local healthcare systems, the problem is not only about model architecture but also about the lack of labeled corpora, the mixing of Vietnamese--English--Latin terminology, and the absence of a local medical knowledge graph --- for a 2024 review to leave no room for this aspect is, in my opinion, a significant gap. Fourth, the paper does not mention Explainable AI (XAI) or uncertainty quantification, even though these two issues are central to modern CDSS systems, especially after the EU AI Act classified clinical decision support systems as high-risk; techniques such as Grad-CAM, conformal prediction, and Dempster-Shafer theory should have been discussed. Fifth, even though this is a review paper, there is no section that clearly explains how the authors selected the literature (no PRISMA methodology or inclusion/exclusion criteria), and the paper also does not position itself against existing surveys such as Litjens 2017 or Esteva 2019, so I cannot tell what new contribution this paper provides compared to those. In terms of presentation, I feel that the historical part covering the 1980s--1990s is somewhat long and rambling while the section on CNN mechanisms --- which should be the important part --- is rather short, and the boxes defining AI/ML/DL feel too basic for a Q1 journal. Overall, the paper still has value as an introductory reference for newcomers, but I think the authors should add a section on Transformers and foundation models, expand the scope to segmentation, include a section on medical AI in low-resource languages, and further discuss XAI and uncertainty before it can be formally published.

\end{document}
